\section{Discussion and Limitations}
\label{sec:discussion}

\subsection{Four vision/training axes, one language axis: localizing the ceiling}

Sec.~\ref{sec:abl-bridge}--\ref{sec:abl-rq6} push on four vision- and
training-side levers that could, in principle, close the remaining F1 gap to
ViMoE-VQA on top of the frozen-backbone, 1-tile bridge --- bridge capacity,
visual-compute allocation, training target and representation alignment --- plus
one language-side lever, decoder capacity. The four vision/training axes are
negative and only the decoder axis is positive --- and even there, only when the
LoRA is placed on the decoder's \emph{attention} projections. The shape of that
split is itself the finding:

\begin{enumerate}
\item \textbf{A bigger bridge does not help; the bridge is not the constraint
  (Sec.~\ref{sec:abl-bridge}).} A $10\times$-larger connector (the 69\,M-parameter
  Full Q-Former) is \emph{worse} than the 7\,M-parameter multi-token bridge on
  both F1 ($-2.2$) and CIDEr-D, and the multi-token bridge already has the lowest
  validation cross-entropy of all five architectures. Whatever the decoder cannot
  do with eight pooled tokens, it also cannot do with sixteen learned queries and
  four fusion layers.
\item \textbf{Reasoning type does not predict visual-compute demand
  (Sec.~\ref{sec:abl-routing}).} Per category, the effect of \ntiles{} on answer
  quality is not significant in any of the eight categories; no learned policy ---
  reasoning-type-informed or not --- beats a fixed \bridge{multi\_token$|$t1} on
  held-out test. This holds on both the original and the tile-count-augmented
  retrained checkpoints, so it is not an artefact of the bridge never having seen
  $>1$ tile during training.
\item \textbf{Multi-reference training and projector-alignment KD do not lift F1
  (Sec.~\ref{sec:abl-rq5}).} Neither training on a resampled reference each epoch
  (\dF{}~$-1.5$) nor distilling the bridge toward Vintern's own pre-aligned
  \texttt{mlp1} projector (\dF{}~$-0.03$, an \emph{absolute} null) moves F1
  upward. The bridge is already close to CE-optimal
  (Sec.~\ref{sec:abl-bridge}) --- there is little room for a training-signal or
  alignment tweak to improve on.
\item \textbf{Decoder-LoRA is the one intervention that moves F1, it is
  bridge-agnostic, and it works only on the attention projections
  (Sec.~\ref{sec:abl-rq6}).} Adapting $0.23\%$ of Qwen2-0.5B's parameters lifts
  F1 $+3.6$ on \multitoken{} and $+5.9$--$7.8$ on the four other bridges --- a
  \emph{larger} effect on the weaker bridges, which rules out ``one bridge's LoRA
  got lucky''. After LoRA all five bridges collapse into a $0.6$-point F1 band and
  a $2.6$-point CIDEr-D band, from plain spreads of ${\sim}4.4$ F1 / ${\sim}13$
  CIDEr-D. Moving the same LoRA budget to the decoder's feed-forward layers
  \emph{diverges} training (val loss $2$--$4$ vs.\ $1.37$, F1 $20$--$38$); the
  useful headroom is specifically in attention, not the decoder broadly. (We flag
  the MLP result as possibly a hyperparameter artefact --- $\alpha=32$ is
  aggressive for the larger intermediate dimension --- so the claim is scoped to
  the recipe's settings.)
\end{enumerate}

\paragraph{Reading.}
Every vision- and training-side lever produces no lift (alignment is a flat
zero); the one lever that touches the decoder produces a lift on every bridge,
and only through its attention projections. That pattern is more informative than
any single ablation: it is not that we tried one clever trick and it happened to
work, it is that \emph{only} the trick which adds attention capacity to the
decoder worked, regardless of which bridge. For this VLM class --- frozen ViT,
frozen small (0.5B) decoder, a few pooled vision tokens --- the frozen decoder is
the ceiling on token-level phrasing match, not the vision pipeline. We report
LoRA as a reference point quantifying that ceiling, not as a replacement for the
paper's main spine (Sec.~\ref{sec:results-main}--\ref{sec:results-corpus}) --- the
frozen, $0.78\%$-param bridge remains the primary contribution, and this finding
explains why that architecture class tops out where it does rather than arguing
it should be abandoned.

\subsection{Relation to ViMoE-VQA's ``reasoning-aware'' claim}

ViMoE-VQA attributes part of its MoE gain to ``approximate reasoning-aware expert
selection''. Our result does not contradict ViMoE's accuracy numbers, but it does
suggest the mechanism is mis-attributed: on the same benchmark, reasoning type
carries no signal about the \emph{optimal} visual-compute action, and ViMoE's own
leave-one-out ablation shows its experts are not strongly specialised. A more
parsimonious account of the MoE gain is added capacity / ensembling, not
reasoning-type routing. Testing this on ViMoE directly (does expert activation
correlate with question type?) is the obvious follow-up and requires only its
router logs.

\subsection{What is robust}

\begin{itemize}
\item The \textbf{grouped split} closes an image-level leakage path; our bridge
  numbers are essentially unchanged from the random-split table
  ($\le 0.5$ F1 / $\le 3$ CIDEr-D). \emph{Caveat:} one early Residual-bridge run
  used in the first comparison was a training-instability outlier (val CE $2.35$
  vs.\ ${\sim}1.5$--$1.7$ elsewhere); with a sound 3-seed run the Residual numbers
  rise sharply (F1 $36.5{\to}45.6$, CIDEr-D $56.3{\to}81.1$) --- so the ``barely
  changed'' claim holds for the four stable bridges, not for that single broken
  run.
\item The \textbf{multi-token bridge} result --- corpus CIDEr-D $92.3$ / BLEU-4
  $18.9$ / ROUGE-L $48.9$ on val, above ViMoE-VQA on all three, at $0.78\%$
  trainable parameters --- holds across 4 seeds and \textbf{on held-out test}
  (F1 $49.2$ vs.\ val $49.6$), a clean, reproducible positive with no val-set
  overfitting.
\item The \textbf{compute-efficiency characterisation}
  (Sec.~\ref{sec:abl-compute}) is the FLOPs/latency analysis ViMoE-VQA explicitly
  deferred.
\item The \textbf{null result on adaptive visual-compute routing}
  (Sec.~\ref{sec:abl-routing}) is cross-validated on two independently-trained
  checkpoint generations --- the same conclusion both times rules out ``the
  bridge just never learned to use tiles'' as a confound.
\item The \textbf{decoder-LoRA result} (Sec.~\ref{sec:abl-rq6}) is the
  most-replicated finding in the paper: $+3.6$ F1 on \multitoken{} (3/3 seeds) and
  $+5.9$--$7.8$ on four other bridges, all converging to a $0.6$-point F1 band;
  confirmed at corpus level and on held-out test. Three checks (seed, bridge,
  metric implementation) plus the attention-vs-MLP contrast all agree.
\end{itemize}

\subsection{Limitations}

\begin{enumerate}
\item \textbf{3-seed, not 5.} Bridge and negative-axis runs are 3 seeds
  (\multitoken{} 4); routing/oracle runs are seed 42. A 5-seed protocol is the
  camera-ready target.
\item \textbf{Frozen 0.5\,B decoder.} The negative routing/training/alignment
  results may not generalise to VLMs with a trainable or larger decoder that
  \emph{can} use extra visual tokens --- and Sec.~\ref{sec:abl-rq6}'s LoRA result
  is direct evidence for exactly that. We claim the negatives for the
  frozen-backbone / small-bridge regime only.
\item \textbf{Decoder-LoRA epoch and rank grid is coarse.} The MLP / attn+MLP
  divergence uses the recipe's $\alpha=32$ and learning rate --- a lower $\alpha$
  or longer schedule might make MLP LoRA viable. Treat the \emph{direction and
  bridge-agnosticism} of the effect as solid, the exact magnitude as provisional.
\item \textbf{$\ntiles \in \{1,3,6\}$ and three oracle bridges.} A finer action
  grid, or bridges that consume per-tile tokens without pooling, might expose
  structure this study cannot see.
\item \textbf{$M(a;x) =$ per-sample CIDEr} is noisy for 4-word answers; a less
  noisy per-sample quality signal (\eg an LLM judge) could in principle recover
  some headroom.
\item \textbf{Cross-paper metric implementations.} Baseline rows are as-reported
  and may use different implementations, especially for METEOR and BLEU; we
  standardise our own numbers on pycocoevalcap.
\item \textbf{Oracle-sweep subset} is equal-per-category; overall numbers are
  re-weighted but tail categories are over-represented in the raw sweep.
\item \textbf{Human validation is done only in reduced form.}
  Sec.~\ref{sec:human} substitutes a single-rater self-check ($N=120$, no image
  access) for the planned 2-annotator, Cohen's-$\kappa$ study; its finding is
  informative but rests on one rater's judgment without an image.
\end{enumerate}

\subsection{Ethical and reproducibility notes}

All models are frozen public checkpoints; only small bridges/heads are trained.
The split script, action space, and oracle tables are released. No human-subjects
data beyond the public AutoViVQA benchmark.
